\heading{第6章-对称性和守恒律}
\begin{question}
在三维空间中考虑宇称算符 $\Pi$。
\begin{enumerate}
\item 证明 $\Pi\psi(\mathbf r)=\psi(-\mathbf r)$ 等同于一次镜面反射再加一次旋转。
\item 写出宇称在球坐标 $(r,\theta,\phi)$ 下的作用。
\item 证明氢原子轨道 $\psi_{n\ell m}=R_{n\ell}(r)Y_\ell^m(\theta,\phi)$ 是宇称本征态，本征值为 $(-1)^\ell$。
\end{enumerate}
\end{question}
\begin{solution}
宇称作用为 $(\boldsymbol r	o-\boldsymbol r)$。它可分解为一次镜面反射和一次绕垂直轴转过 $\pi$ 的旋转；例如先令 $z\to-z$，再绕 $z$ 轴转 $\pi$，总体效果就是 $(x,y,z)\to(-x,-y,-z)$。

球坐标中
\[
(r,\theta,\phi)\longmapsto (r,\pi-\theta,\phi+\pi),
\]
所以
\[
\Pi\psi(r,\theta,\phi)=\psi(r,\pi-\theta,\phi+\pi).
\]
又
\[
Y_\ell^m(\pi-\theta,\phi+\pi)=(-1)^\ell Y_\ell^m(\theta,\phi),
\]
因此中心势定态 $R_{n\ell}(r)Y_\ell^m(\theta,\phi)$ 是宇称本征态，宇称本征值为 $(-1)^\ell$。
\end{solution}

\begin{question}
设 $Q$ 是厄米算符，证明
\[
U=\exp(\ii Q)
\]
是幺正算符。提示：先由幂级数证明 $U^\dagger=\exp(-\ii Q)$。
\end{question}
\begin{solution}
若 $Q=Q^\dagger$，令 $U=\exp(\ii Q)$。由幂级数取伴随得
\[
U^\dagger=\left(\sum_{n=0}^\infty { (\ii Q)^n\over n!}\right)^\dagger
=\sum_{n=0}^\infty { (-\ii Q)^n\over n!}=\exp(-\ii Q).
\]
因为 $Q$ 与自身任意函数对易，
\[
U^\dagger U=\exp(-\ii Q)\exp(\ii Q)=1,
\qquad UU^\dagger=1.
\]
所以 $U$ 是么正算符。
\end{solution}

\begin{question}
证明动量算符在空间平移下不变：若 $T(a)=\exp(-\ii ap/\hbar)$，则
\[
T(a)^\dagger pT(a)=p.
\]
\end{question}
\begin{solution}
平移算符 $T(a)=\exp(-\ii ap/\hbar)$ 是动量 $p$ 的函数，故 $[p,T(a)]=0$。于是
\[
p'=T(a)^\dagger pT(a)=pT(a)^\dagger T(a)=p.
\]
这说明空间平移不改变动量算符本身，只改变位置算符的原点。
\end{solution}

\begin{question}
设算符 $Q(x,p)$ 可展开为 $x$ 与 $p$ 的幂级数，证明
\[
T(a)^\dagger Q(x,p)T(a)=Q(x+a,p).
\]
并由此说明一维哈密顿量 $H=p^2/(2m)+V(x)$ 在平移下如何变换。
\end{question}
\begin{solution}
由 $T(a)^\dagger xT(a)=x+a$ 与 $T(a)^\dagger pT(a)=p$，对任意可展开为幂级数的算符函数
\[
Q(x,p)=\sum_{m,n}c_{mn}x^mp^n
\]
逐项变换：
\[
T^\dagger Q(x,p)T=\sum_{m,n}c_{mn}(x+a)^mp^n=Q(x+a,p).
\]
因此
\[
H'={p^2\over 2m}+V(x+a),
\]
这就是式 (6.8) 的内容；若 $V(x+a)=V(x)$，则 $H'=H$，即 $[H,T(a)]=0$。
\end{solution}

\begin{question}
由布洛赫条件
\[
\psi(x+a)=e^{\ii qa}\psi(x)
\]
证明波函数可写为
\[
\psi(x)=e^{\ii qx}u(x),\qquad u(x+a)=u(x).
\]
\end{question}
\begin{solution}
由式 (6.11) $\psi(x+a)=e^{\ii qa}\psi(x)$，定义
\[
u(x)=e^{-\ii qx}\psi(x).
\]
则
\[
u(x+a)=e^{-\ii q(x+a)}\psi(x+a)=e^{-\ii qx}e^{-\ii qa}e^{\ii qa}\psi(x)=u(x).
\]
所以 $u$ 与势具有相同周期，且
\[
\psi(x)=e^{\ii qx}u(x),
\]
即布洛赫形式。反过来若 $u(x+a)=u(x)$，立刻得到 $\psi(x+a)=e^{\ii qa}\psi(x)$。
\end{solution}

\begin{question}
质量为 $m$ 的粒子在周期为 $a$ 的一维势场 $V(x)$ 中运动。设
\[
\psi_{nq}(x)=e^{\ii qx}u_{nq}(x),\qquad u_{nq}(x+a)=u_{nq}(x).
\]
\begin{enumerate}
\item 推导 $u_{nq}$ 满足的本征值方程。
\item 说明如何用有限差分方法在一个周期内数值求解能带 $E_{nq}$。
\item 对不同 $q$ 值画出最低能带的示意曲线。
\end{enumerate}
\end{question}
\begin{solution}
把 $\psi_{nq}=e^{\ii qx}u_{nq}$ 代入薛定谔方程，得到作用在周期函数 $u_{nq}$ 上的本征值方程
\[
\left[-{\hbar^2\over 2m}\left({\dd\over\dd x}+\ii q\right)^2+V(x)\right]u_{nq}=E_{nq}u_{nq}.
\]
展开后即
\[
-{\hbar^2\over 2m}u''-{\ii\hbar^2q\over m}u'+{\hbar^2q^2\over 2m}u+Vu=Eu.
\]
数值上取一个周期 $[0,a]$，令 $u(0)=u(a)$，用周期边界条件构造差分矩阵：
\[
u'_j\simeq {u_{j+1}-u_{j-1}\over 2\Delta x},\qquad
u''_j\simeq {u_{j+1}-2u_j+u_{j-1}\over \Delta x^2},
\]
首尾指标按模 $N$ 处理。对每个 $q$ 对角化该矩阵，最低两个本征值就是 $E_1(q),E_2(q)$。由于 $q$ 与 $q+2\pi/a$ 等价，只需画第一布里渊区 $-\pi/a\le q\le\pi/a$。
\end{solution}

\begin{question}
考虑一维中两个质量分别为 $m_1,m_2$ 的粒子，哈密顿量中相互作用只依赖 $|x_1-x_2|$。
\begin{enumerate}
\item 求同时平移两个粒子的平移算符，并指出其生成元。
\item 证明该哈密顿量与总动量对易。
\item 说明这对应于系统的整体平移不变性。
\end{enumerate}
\end{question}
\begin{solution}
两粒子同时平移 $a$：
\[
T(a)\psi(x_1,x_2)=\psi(x_1-a,x_2-a).
\]
对无穷小 $a$ 展开，
\[
T(a)=1-a(\partial_{x_1}+\partial_{x_2})+O(a^2)
=1-{\ii a\over\hbar}(p_1+p_2)+O(a^2),
\]
故
\[
T(a)=\exp[-\ii a(p_1+p_2)/\hbar].
\]
哈密顿量中的动能显然与 $P=p_1+p_2$ 对易，而势能只依赖 $x_1-x_2$，也有
\[
[P,V(|x_1-x_2|)]=-\ii\hbar(\partial_{x_1}+\partial_{x_2})V=0.
\]
所以 $[H,P]=0$，总动量守恒。
\end{solution}

\begin{question}
\begin{enumerate}
\item 证明宇称算符 $\Pi$ 是厄米算符。
\item 证明宇称算符的本征值只能为 $+1$ 或 $-1$。
\end{enumerate}
\end{question}
\begin{solution}
宇称算符满足 $\Pi^2=1$ 且保持内积：
\[
\langle \Pi f|\Pi g\rangle=\int f^*(-x)g(-x)\,\dd x=\int f^*(x)g(x)\,\dd x.
\]
故 $\Pi$ 是么正的。又 $\Pi^{-1}=\Pi$，于是 $\Pi^\dagger=\Pi$，它也是厄米的。

若 $\Pi\psi=\lambda\psi$，再作用一次 $\Pi^2\psi=\lambda^2\psi=\psi$，所以
\[
\lambda=\pm1.
\]
\end{solution}

\begin{question}
真标量在宇称下不变，赝标量在宇称下变号。证明：真标量 $f$ 满足 $[\Pi,f]=0$，赝标量 $g$ 满足 $\{\Pi,g\}=0$。并讨论真矢量与赝矢量在宇称下的变换。
\end{question}
\begin{solution}
若 $f$ 是真标量，宇称下不变：$\Pi^\dagger f\Pi=f$。因 $\Pi^\dagger=\Pi$ 且 $\Pi^2=1$，得 $\Pi f=f\Pi$，即 $[\Pi,f]=0$。

若 $f$ 是赝标量，宇称下变号：$\Pi^\dagger f\Pi=-f$，等价于 $\Pi f=-f\Pi$，即 $\{\Pi,f\}=0$。

真矢量每个分量在宇称下变号，故与 $\Pi$ 反对易；赝矢量在宇称下不变，故与 $\Pi$ 对易。
\end{solution}

\begin{question}
证明位置算符和动量算符都具有奇宇称，即在一维中
\[
\Pi x\Pi=-x,
\qquad
\Pi p\Pi=-p,
\]
并把结论推广到三维情形。
\end{question}
\begin{solution}
在一维中
\[
(\Pi x\Pi\psi)(x)=\Pi x\psi(-x)=\Pi[-x\psi(-x)]=-x\psi(x),
\]
所以 $\Pi x\Pi=-x$。又 $p=-\ii\hbar\partial_x$，
\[
(\Pi p\Pi\psi)(x)=\Pi[-\ii\hbar\partial_x\psi(-x)]=\ii\hbar\psi'(x)=-p\psi(x),
\]
故 $\Pi p\Pi=-p$。三维逐分量完全相同：
\[
\Pi\mathbf r\Pi=-\mathbf r,
\qquad
\Pi\mathbf p\Pi=-\mathbf p.
\]
\end{solution}

\begin{question}
考虑确定宇称态 $|i\rangle,|f\rangle$ 之间的矩阵元 $\langle f|Q|i\rangle$。根据算符 $Q$ 的宇称，推导该矩阵元非零的选择定则。说明该规则如何用于真标量、真矢量和赝矢量算符。
\end{question}
\begin{solution}
设初、末态宇称分别为 $\eta_i,\eta_f=\pm1$，算符 $Q$ 的宇称为 $\eta_Q$。矩阵元
\[
M=\langle f|Q|i\rangle
\]
在插入 $\Pi^2=1$ 后变为
\[
M=\eta_f\eta_Q\eta_i M.
\]
因此非零的必要条件是
\[
\eta_f\eta_Q\eta_i=1.
\]
真标量 $\eta_Q=+1$，要求初末态同宇称；真矢量 $\eta_Q=-1$，要求初末态异宇称。对中心势态，这就是 $(-1)^{\ell'+\ell}\eta_Q=1$。
\end{solution}

\begin{question}
自旋角动量与宇称算符对易：$[\Pi,\mathbf S]=0$。在自旋 $1/2$ 的标准基中，说明宇称对二分量旋量的作用，并证明该作用可取为单位矩阵乘以一个整体常数。
\end{question}
\begin{solution}
宇称只改变空间坐标，不改变内禀自旋自由度。因此
\[
[\Pi,S_i]=0\qquad (i=x,y,z).
\]
在 $S_z$ 的标准基 $\{|+\rangle,|-\rangle\}$ 中，$\Pi$ 与所有泡利矩阵对易。与全部 $2\times2$ 自旋算符都对易的矩阵只能是单位阵的倍数。约定把这个倍数取为 $+1$，于是自旋波函数的宇称作用就是恒等作用。
\end{solution}

\begin{question}
考虑氢原子中的电子。
\begin{enumerate}
\item 若电子处于基态，不做显式积分证明 $\langle z\rangle=0$。
\item 说明当电子处于 $n=2$ 简并子空间时，$\langle z\rangle$ 不一定为零；给出一个例子并计算其值。
\end{enumerate}
\end{question}
\begin{solution}
(a) 氢原子基态是 $\ell=0$ 的偶宇称态，而 $z$ 是奇宇称算符，所以
\[
\langle z\rangle=0.
\]

(b) $n=2$ 能级中 $2s$ 与 $2p$ 简并，可以取没有确定宇称的线性组合。例取
\[
\psi={1\over\sqrt2}(\psi_{200}+\psi_{210}).
\]
则 $\langle p_z\rangle=0$ 仍成立，因为 $p_z$ 只联结角动量相差 $1$ 的态而径向与角向相位使本例为实驻波。若题意为位置偶极矩，则
\[
\langle z\rangle={1\over2}\big(\langle200|z|210\rangle+\langle210|z|200\rangle\big)=-3a,
\]
其中 $a$ 为玻尔半径。这个例子说明简并子空间内可构造非定宇称态，其奇算符期望值不必为零。
\end{solution}

\begin{question}
考察二维平面中绕原点的无穷小旋转矩阵。对角化该矩阵，并把它与角动量本征态在无穷小旋转下获得相位的描述联系起来。
\end{question}
\begin{solution}
矩阵
\[
M=\begin{pmatrix}1&-\epsilon\\ \epsilon&1\end{pmatrix}
\]
的本征值为 $1\pm \ii\epsilon$，本征矢可取 $2^{-1/2}(1,\mp\ii)^T$。因此
\[
M=S\begin{pmatrix}1+\ii\epsilon&0\\0&1-\ii\epsilon\end{pmatrix}S^{-1}.
\]
若把有限转角写作 $N\epsilon=\phi$，则
\[
\lim_{N\to\infty}(1\pm \ii\phi/N)^N=e^{\pm\ii\phi}.
\]
变回原基得
\[
\lim_{N\to\infty}M^N=\begin{pmatrix}\cos\phi&-\sin\phi\\ \sin\phi&\cos\phi\end{pmatrix},
\]
即有限旋转矩阵。
\end{solution}

\begin{question}
证明标量算符在任意无穷小旋转下保持不变，当且仅当它与角动量各分量对易：
\[
[L_i,f]=0\qquad (i=x,y,z).
\]
\end{question}
\begin{solution}
标量算符在旋转下满足
\[
U(R)^\dagger fU(R)=f.
\]
对无穷小旋转 $U=1-\ii\delta\boldsymbol\phi\cdot\mathbf L/\hbar$，保留一阶项：
\[
f+{\ii\over\hbar}\delta\boldsymbol\phi\cdot[\mathbf L,f]=f.
\]
任意 $\delta\boldsymbol\phi$ 下都成立，所以
\[
[L_i,f]=0\quad(i=x,y,z).
\]
反过来若这些对易子全为零，指数形式的有限旋转也使 $f$ 不变。
\end{solution}

\begin{question}
从矢量算符的定义
\[
[L_i,V_j]=\ii\hbar\epsilon_{ijk}V_k
\]
出发，求矢量算符绕 $y$ 轴作无穷小旋转时各分量的变化。
\end{question}
\begin{solution}
绕 $y$ 轴转无穷小角 $\delta\phi$ 时，真矢量满足经典旋转律
\[
V_x' = V_x+\delta\phi V_z,
\quad V_y'=V_y,
\quad V_z'=V_z-\delta\phi V_x.
\]
量子形式为
\[
V_i'=V_i+{\ii\over\hbar}\delta\phi [L_y,V_i].
\]
由矢量算符定义 $[L_i,V_j]=\ii\hbar\epsilon_{ijk}V_k$，立即得到上述三个分量变换。
\end{solution}

\begin{question}
考虑绕 $y$ 轴的无穷小旋转作用在角动量本征态 $|\ell m\rangle$ 上。利用 $L_y=(L_+-L_-)/(2\ii)$，证明旋转只把 $m$ 与 $m\pm1$ 的态混合，并写出一阶表达式。
\end{question}
\begin{solution}
绕 $y$ 轴无穷小旋转
\[
R_y(\epsilon)=1-{\ii\epsilon\over\hbar}L_y.
\]
用 $L_y=(L_+-L_-)/(2\ii)$ 及
\[
L_\pm|\ell m\rangle=\hbar\sqrt{\ell(\ell+1)-m(m\pm1)}|\ell,m\pm1\rangle,
\]
得
\[
R_y(\epsilon)|n\ell m\rangle=
|n\ell m\rangle-{\epsilon\over2}
\sqrt{\ell(\ell+1)-m(m+1)}|n\ell,m+1\rangle
\]
\[
\hspace{5em}+{\epsilon\over2}
\sqrt{\ell(\ell+1)-m(m-1)}|n\ell,m-1\rangle+O(\epsilon^2).
\]
因此 $D$ 矩阵只在 $m'=m,m\pm1$ 处有一阶非零元，且 $n,\ell$ 不变。
\end{solution}

\begin{question}
一维自由粒子哈密顿量 $H=p^2/(2m)$ 同时具有平移和宇称对称性。
\begin{enumerate}
\item 证明平移算符与宇称算符一般不对易。
\item 从动量本征态和宇称本征态两种角度解释自由粒子能谱的二重简并。
\end{enumerate}
\end{question}
\begin{solution}
平移 $T(a)$ 与宇称 $\Pi$ 满足
\[
\Pi T(a)\Pi=T(-a),
\]
所以一般不对易。自由粒子哈密顿量 $H=p^2/2m$ 同时与二者对易。

动量本征态 $|p\rangle$ 与 $|-p\rangle$ 能量相同：$E=p^2/2m$，宇称正是把二者互换。另一方面，取宇称本征态
\[
\cos(px/\hbar),\qquad \sin(px/\hbar),
\]
平移会把它们线性混合。由于对称操作把某个本征态变为同能量的不同本征态，简并性随之出现。
\end{solution}

\begin{question}
对任意矢量算符 $\mathbf V$ 定义
\[
V_\pm=V_x\pm\ii V_y.
\]
\begin{enumerate}
\item 利用 $[L_i,V_j]=\ii\hbar\epsilon_{ijk}V_k$ 推导 $V_\pm,V_z$ 与 $L_z,L^2$ 的有关对易关系。
\item 说明 $V_q$ 对角动量量子数 $m$ 与 $\ell$ 的选择定则。
\end{enumerate}
\end{question}
\begin{solution}
定义 $V_\pm=V_x\pm\ii V_y$。由 $[L_i,V_j]=\ii\hbar\epsilon_{ijk}V_k$ 得
\[
[L_z,V_\pm]=\pm\hbar V_\pm,
\qquad [L_z,V_z]=0.
\]
继续计算可得 $V_\pm$ 改变量子数 $m$ 为 $m\pm1$，$V_z$ 保持 $m$ 不变。再由 $[L^2,V_q]$ 的关系知矢量算符只能把 $\ell$ 变为 $\ell,\ell\pm1$；结合宇称时，对位置这样的真矢量还要求 $\ell+\ell'$ 为奇数，故通常为 $\Delta\ell=\pm1$。
\end{solution}

\begin{question}
证明若标量算符 $f$ 满足 $[L_z,f]=0$ 和 $[L^2,f]=0$，则其在角动量本征态之间的矩阵元只能在 $m'=m$ 且 $\ell'=\ell$ 时非零。
\end{question}
\begin{solution}
若 $f$ 为标量，$[L_z,f]=0$ 给出
\[
0=\langle \ell'm'|[L_z,f]|\ell m\rangle
=\hbar(m'-m)\langle \ell'm'|f|\ell m\rangle,
\]
所以非零时 $m'=m$。又 $[L^2,f]=0$ 给出
\[
\hbar^2[\ell'(\ell'+1)-\ell(\ell+1)]\langle \ell'm'|f|\ell m\rangle=0,
\]
所以 $\ell'=\ell$。这就是标量算符的选择定则。
\end{solution}

\begin{question}
电子处于氢原子态
\[
\psi={1\over\sqrt2}\left(\psi_{211}+\psi_{21,-1}\right).
\]
利用选择定则或直接计算，求 $\langle\mathbf r\rangle$。
\end{question}
\begin{solution}
位置 $\mathbf r$ 是真矢量、奇宇称算符。给出的态只含有 $\ell=1$ 的氢原子态，任意 $\langle n\ell m|\mathbf r|n\ell m'\rangle$ 中初末态宇称同为 $(-1)^\ell=-1$，乘上奇算符后整体为奇函数，积分为零。因此
\[
\langle \mathbf r\rangle=0.
\]
这也可由维格纳--埃卡特定理看出：位置算符要求 $\Delta\ell=\pm1$，而该线性组合内部只有 $\ell=1$ 分量。
\end{solution}

\begin{question}
\begin{enumerate}
\item 由矢量算符定义推导 $V_z,V_\pm$ 与 $L_z,L_\pm$ 的六个对易关系。
\item 由这些对易关系推导矢量算符矩阵元的选择定则。
\end{enumerate}
\end{question}
\begin{solution}
矢量算符定义给出
\[
[L_z,V_z]=0,
\quad [L_z,V_\pm]=\pm\hbar V_\pm,
\]
\[
[L_\pm,V_z]=\mp\hbar V_\pm,
\quad [L_\pm,V_\mp]=\pm2\hbar V_z,
\quad [L_\pm,V_\pm]=0.
\]
这些式子正是球基分量 $V_q(q=0,\pm1)$ 的角动量张量变换律。进一步把它们夹在 $|\ell m\rangle$ 态之间，即得到式 (6.57) 以及相应的选择定则。
\end{solution}

\begin{question}
把两个角动量 $j_1,j_2$ 耦合成总角动量 $J$：
\[
|JM\rangle=\sum_{m_1,m_2} C^{JM}_{j_1m_1j_2m_2}|j_1m_1\rangle|j_2m_2\rangle.
\]
证明 Clebsch--Gordan 系数非零时必须满足 $M=m_1+m_2$，并推导由升降算符给出的递推关系。
\end{question}
\begin{solution}
把
\[
|JM\rangle=\sum_{m_1,m_2}C^{JM}_{j_1m_1j_2m_2}|j_1m_1\rangle|j_2m_2\rangle
\]
代入 $J_z=J_{1z}+J_{2z}$，得
\[
(m_1+m_2)C^{JM}_{j_1m_1j_2m_2}=M C^{JM}_{j_1m_1j_2m_2},
\]
所以非零 CG 系数满足 $m_1+m_2=M$。再将 $J_\pm=J_{1\pm}+J_{2\pm}$ 作用到上式两边，并比较 $|j_1m_1\rangle|j_2m_2\rangle$ 的系数，即得标准递推关系。该递推关系和归一化条件共同确定 CG 系数。
\end{solution}

\begin{question}
将矢量算符球分量的对易关系夹在角动量本征态之间，推导其矩阵元之间的关系，并说明这些关系如何导向维格纳--埃卡特定理的形式。
\end{question}
\begin{solution}
将六个对易关系夹在 $\langle n'\ell'm'|$ 与 $|n\ell m\rangle$ 之间，左边用 $L_z,L_\pm$ 对角或升降作用，右边保留矢量算符矩阵元。例如
\[
\hbar(m'-m)\langle n'\ell'm'|V_+|n\ell m\rangle
=\hbar\langle n'\ell'm'|V_+|n\ell m\rangle,
\]
说明 $V_+$ 只允许 $m'=m+1$。其余五式同理给出 $V_0,V_\pm$ 的全部 $m$ 选择规则和相邻矩阵元比例。与习题 6.23 的 CG 递推关系比较，即得维格纳--埃卡特定理中的角向因子。
\end{solution}

\begin{question}
把位置算符 $\mathbf r$ 写成球张量分量 $r_0,r_\pm$，并利用维格纳--埃卡特定理推导氢原子电偶极矩阵元的选择定则。
\end{question}
\begin{solution}
位置算符的球分量为
\[
r_0=z=r\sqrt{4\pi/3}\,Y_1^0,
\qquad
r_\pm=\mp {x\pm\ii y\over\sqrt2}=r\sqrt{4\pi/3}\,Y_1^{\pm1}.
\]
因此它是 rank-$1$ 球张量。维格纳--埃卡特定理给出
\[
\langle n'\ell'm'|r_q|n\ell m\rangle
=\langle \ell m;1q|\ell'm'\rangle\langle n'\ell'\|r\|n\ell\rangle.
\]
选择定则为 $m'=m+q$、$\ell'=\ell\pm1$，并由宇称排除 $\ell'=\ell$。具体径向数值都包含在约化矩阵元中。
\end{solution}

\begin{question}
求例题 6.7 中系统的海森伯绘景动量算符 $p_H(t)$，并讨论所得结果与经典运动方程的对应。
\end{question}
\begin{solution}
海森伯绘景中
\[
{\dd p_H\over\dd t}={\ii\over\hbar}[H,p_H].
\]
若例题中的哈密顿量为 $H=p^2/2m+V(x)$，则
\[
{\dd p_H\over\dd t}=-V'(x_H).
\]
对线性势 $V=-Fx$，这给出 $p_H(t)=p(0)+Ft$；对谐振子 $V=m\omega^2x^2/2$，得到 $\dot p_H=-m\omega^2x_H$。两者都与经典哈密顿方程同形。
\end{solution}

\begin{question}
考虑质量为 $m$ 的一维自由粒子。证明海森伯绘景中
\[
x_H(t)=x_H(0)+{p_H(0)\over m}t,
\qquad
p_H(t)=p_H(0),
\]
并解释它们与经典自由运动方程的关系。
\end{question}
\begin{solution}
自由粒子 $H=p^2/2m$。海森伯方程给
\[
\dot p_H={\ii\over\hbar}[H,p_H]=0,
\qquad
\dot x_H={\ii\over\hbar}[H,x_H]={p_H\over m}.
\]
积分得
\[
p_H(t)=p_H(0),
\qquad
x_H(t)=x_H(0)+{p_H(0)\over m}t.
\]
这与经典自由粒子的 $p=$ 常量、$x=x_0+pt/m$ 完全一致。
\end{solution}

\begin{question}
由薛定谔方程对 $\psi(x,t+\delta t)$ 作泰勒展开，证明无穷小时间平移算符为
\[
U(\delta t)=1-{\ii H\delta t\over\hbar}+O(\delta t^2).
\]
\end{question}
\begin{solution}
薛定谔绘景中
\[
\psi(t+\delta t)=\psi(t)+\delta t\,{\partial\psi\over\partial t}+O(\delta t^2)
=\left(1-{\ii H\delta t\over\hbar}\right)\psi(t)+O(\delta t^2).
\]
所以无穷小时间平移算符为
\[
U(\delta t)=1-{\ii H\delta t\over\hbar},
\]
有限时间演化由连续相乘给出
\[
U(t)=\exp(-\ii Ht/\hbar)
\]
（当 $H$ 不显含时时）。
\end{solution}

\begin{question}
对海森伯绘景算符 $Q_H(t)=U^\dagger(t)QU(t)$ 求导，推导海森伯运动方程。并将其用于 $H=p^2/(2m)+V(x)$，求 $x_H,p_H$ 的运动方程。
\end{question}
\begin{solution}
由
\[
Q_H(t)=U^\dagger(t)QU(t)
\]
求导，且 $\dot U=-\ii HU/\hbar$，得
\[
{\dd Q_H\over\dd t}={\ii\over\hbar}[H,Q_H]+\left({\partial Q\over\partial t}\right)_H.
\]
对 $H=p^2/2m+V(x)$，
\[
\dot x_H={p_H\over m},
\qquad
\dot p_H=-\left({\partial V\over\partial x}\right)_{x=x_H}.
\]
这就是算符形式的牛顿方程。
\end{solution}

\begin{question}
设一维不含时哈密顿量的定态为 $\psi_n(x)$、能量为 $E_n$。
\begin{enumerate}
\item 证明传播子可写为
\[
K(x,x';t)=\sum_n\psi_n(x)\psi_n^*(x')e^{-\ii E_nt/\hbar}.
\]
\item 将该表达式用于一维谐振子，得到标准传播子的形式。
\end{enumerate}
\end{question}
\begin{solution}
若 $H\psi_n=E_n\psi_n$ 且 $\{\psi_n\}$ 完备，则
\[
\psi(x,t)=\int K(x,x',t)\psi(x',0)\,\dd x',
\]
其中
\[
K(x,x',t)=\sum_n\psi_n(x)\psi_n^*(x')e^{-\ii E_nt/\hbar}.
\]
谐振子的 Mehler 求和给出标准传播子
\[
K=\sqrt{\frac{m\omega}{2\pi\ii\hbar\sin\omega t}}
\exp\left\{\frac{\ii m\omega}{2\hbar\sin\omega t}
\big[(x^2+x'^2)\cos\omega t-2xx'\big]\right\}.
\]
它在 $t\to0$ 时趋于 $\delta(x-x')$。
\end{solution}

\begin{question}
用动量完备基定义平移算符
\[
T(a)=\int e^{-\ii pa/\hbar}|p\rangle\langle p|\,\dd p.
\]
证明它对任意波函数的作用为 $T(a)\psi(x)=\psi(x-a)$。
\end{question}
\begin{solution}
用动量完备性定义
\[
T(a)=\int e^{-\ii pa/\hbar}|p\rangle\langle p|\,\dd p.
\]
若
\[
\psi(x)=\int \phi(p){e^{\ii px/\hbar}\over\sqrt{2\pi\hbar}}\,\dd p,
\]
则
\[
T(a)\psi(x)=\int \phi(p){e^{\ii p(x-a)/\hbar}\over\sqrt{2\pi\hbar}}\,\dd p=\psi(x-a).
\]
这个推导只要求傅里叶展开存在，不要求 $\psi$ 处处可微；所以即使在某点一阶导数不连续，谱分解定义仍给出正确平移。
\end{solution}

\begin{question}
自旋 $1/2$ 的旋量在绕单位矢量 $\hat n$ 转动角 $\phi$ 时由
\[
U(\phi)=\exp\left(-{\ii\phi\over2}\boldsymbol\sigma\cdot\hat n\right)
\]
给出。
\begin{enumerate}
\item 证明泡利矩阵恒等式 $(\boldsymbol\sigma\cdot\mathbf a)(\boldsymbol\sigma\cdot\mathbf b)=\mathbf a\cdot\mathbf b+\ii\boldsymbol\sigma\cdot(\mathbf a\times\mathbf b)$。
\item 推导 $U(\phi)$ 的闭式表达式和矩阵形式。
\end{enumerate}
\end{question}
\begin{solution}
泡利矩阵满足
\[
(\boldsymbol\sigma\cdot\mathbf a)(\boldsymbol\sigma\cdot\mathbf b)
=\mathbf a\cdot\mathbf b+\ii\boldsymbol\sigma\cdot(\mathbf a\times\mathbf b).
\]
取 $\mathbf a=\mathbf b=\hat n$ 得 $(\boldsymbol\sigma\cdot\hat n)^2=1$。因此
\[
\exp\left(-{\ii\phi\over2}\boldsymbol\sigma\cdot\hat n\right)
=\cos{\phi\over2}-\ii(\boldsymbol\sigma\cdot\hat n)\sin{\phi\over2}.
\]
若 $\hat n=(\sin\theta\cos\varphi,\sin\theta\sin\varphi,\cos\theta)$，在 $S_z$ 基下
\[
R_{\hat n}(\phi)=\begin{pmatrix}
\cos{\phi\over2}-\ii\cos\theta\sin{\phi\over2}&-\ii e^{-\ii\varphi}\sin\theta\sin{\phi\over2}\\
-\ii e^{\ii\varphi}\sin\theta\sin{\phi\over2}&\cos{\phi\over2}+\ii\cos\theta\sin{\phi\over2}
\end{pmatrix}.
\]
\end{solution}

\begin{question}
考虑边长为 $L$ 的二维无限深方势阱，原点取在方阱中心。
\begin{enumerate}
\item 说明为什么能量只依赖 $n_x^2+n_y^2$，并证明 $(a,b)$ 与 $(b,a)$ 态简并。
\item 证明绕中心旋转 $90^\circ$ 会把这两个态相互变换。
\item 构造旋转算符的好态线性组合。
\end{enumerate}
\end{question}
\begin{solution}
二维方阱的能量只依赖 $n_x^2+n_y^2$，所以 $(a,b)$ 与 $(b,a)$ 简并。绕中心转 $90^\circ$ 把坐标 $(x,y)$ 变到 $(-y,x)$，从而把两个量子数互换；相差的符号由正弦函数奇偶性决定。

取
\[
\psi_\pm={1\over\sqrt2}(\psi_{ab}\pm\psi_{ba}).
\]
若 $a,b$ 同奇或同偶，$90^\circ$ 旋转后这两个组合只乘上相位（本征值为 $\pm1$ 或 $\pm\ii$，取决于边界坐标约定），因此它们是旋转算符的本征态。例 $a=5,b=3$ 时，简并二重态可用这些对称/反对称组合对角化旋转操作。
\end{solution}

\begin{question}
库仑势具有比普通旋转不变性更高的对称性。定义拉普拉斯--龙格--楞次矢量
\[
\mathbf M={1\over2m}(\mathbf p\times\mathbf L-\mathbf L\times\mathbf p)-k{\mathbf r\over r}.
\]
证明在 $V(r)=-k/r$ 中 $\mathbf M$ 与哈密顿量对易，并推导它与 $\mathbf L$ 之间的主要对易关系。
\end{question}
\begin{solution}
拉普拉斯--龙格--楞次矢量
\[
\mathbf M={1\over2m}(\mathbf p\times\mathbf L-\mathbf L\times\mathbf p)-k{\mathbf r\over r}
\]
只在库仑势 $V=-k/r$ 中守恒。直接用
\[
[x_i,p_j]=\ii\hbar\delta_{ij},\qquad [L_i,L_j]=\ii\hbar\epsilon_{ijk}L_k,
\]
以及 $[p_i,f(r)]=-\ii\hbar\partial_i f$ 计算，可得
\[
[H,L_i]=0,
\qquad
[H,M_i]=0,
\]
并且 $\mathbf M$ 像矢量一样在旋转下变换：
\[
[L_i,M_j]=\ii\hbar\epsilon_{ijk}M_k.
\]
进一步整理 $[M_i,M_j]$ 得
\[
[M_i,M_j]=-\frac{2\ii\hbar H}{m}\epsilon_{ijk}L_k.
\]
这些对易关系解释了氢原子能级中比普通转动对称更高的简并度。
\end{solution}

\begin{question}
伽利略变换把参考系变到相对原参考系以速度 $v$ 运动的新参考系。
\begin{enumerate}
\item 求无穷小 boost 下 $x$ 与 $p$ 的变换，并解释其物理意义。
\item 利用 Baker--Campbell--Hausdorff 公式整理有限 boost 算符的形式。
\item 证明若 $\psi$ 满足势 $V(x)$ 下的薛定谔方程，则 boost 后的波函数满足相应移动势场中的薛定谔方程。
\end{enumerate}
\end{question}
\begin{solution}
伽利略 boost 到速度 $v$ 的参考系可写作
\[
G(v,t)=\exp\left[-{\ii\over\hbar}\left(mvx-{1\over2}mv^2t\right)\right]\,T(vt).
\]
对无穷小 $v$，
\[
x'=G^\dagger xG=x-vt,
\qquad
p'=G^\dagger pG=p-mv.
\]
这正是经典伽利略变换：新参考系中位置坐标少了 $vt$，动量少了 $mv$。利用 BCH 公式可把相位算符和平移算符合并为题中给出的形式。
\end{solution}

\begin{question}
讨论时间反演不变性。以无空气阻力下抛体运动为例说明：若在某一时刻把速度反向，粒子会沿原轨迹反向返回。把这一思想推广到量子力学，说明时间反演对 $x,p$ 和虚数单位 $\ii$ 的作用。
\end{question}
\begin{solution}
时间反演把运动方向反过来：
\[
\Theta x\Theta^{-1}=x,
\qquad
\Theta p\Theta^{-1}=-p,
\qquad
\Theta \ii\Theta^{-1}=-\ii.
\]
若势能 $V(x)$ 为实函数且不含速度，哈密顿量 $H=p^2/2m+V(x)$ 在时间反演下不变。于是若 $x(t)$ 是一条经典轨迹，令某时刻速度反号，随后轨迹就是原轨迹的反向重放。量子力学中复共轭的出现正是为了使 $p=-\ii\hbar\partial_x$ 在时间反演下变号。
\end{solution}

\begin{question}
自旋 $1/2$ 粒子的时间反演算符可取为 $\Theta=-\ii\sigma_yK$，其中 $K$ 表示复共轭。
\begin{enumerate}
\item 证明 $\Theta^2=-1$。
\item 若哈密顿量与时间反演对易，证明所有能级至少二重简并（Kramers 简并）。
\end{enumerate}
\end{question}
\begin{solution}
自旋 $1/2$ 的时间反演算符可取
\[
\Theta=-\ii\sigma_y K,
\]
其中 $K$ 为复共轭。于是
\[
\Theta^2=(-\ii\sigma_y K)(-\ii\sigma_y K)=(-\ii\sigma_y)(-\ii\sigma_y)^*=-1.
\]
若 $H$ 与 $\Theta$ 对易，且 $H|\psi\rangle=E|\psi\rangle$，则 $H\Theta|\psi\rangle=E\Theta|\psi\rangle$。若二者相同到一个相位，$\Theta|\psi\rangle=c|\psi\rangle$，再作用一次给 $\Theta^2|\psi\rangle=|c|^2|\psi\rangle$，与 $\Theta^2=-1$ 矛盾。因此 $|\psi\rangle$ 与 $\Theta|\psi\rangle$ 是不同态且同能量，这就是 Kramers 简并。
\end{solution}
